\section{GC-04原生组织与端点计算状态}

本例保留\sid{GC-04}的65 nm current-programmed dynamic-cascode 3T1C gain-cell，采用二进制端点、伪差分ACIM；数字外围使用共同28 nm能力预算。原宏为64行、64个pair列，保留该局部读写负载；八个位平面构成 $K=N=64$ 的INT8逻辑矩阵。有效容量4096 Byte、输入64 Byte，输出64个22-bit容器。八个 $64\times64$ pair tile共32768 pair、65536个3T1C cell，位平面共同编码一份矩阵。

原cell支持0--700 nA八电流级，两cell伪差分形成15个有符号等级（GC-04 p.6 Fig.8）。本例只用端点：bit1存 $(700,0)$ nA，bit0存 $(0,700)$ nA；每INT8权重16cell，信号编码不增加有效payload。原生64行/64列的选择保留测量局部负载，同时避免为统一容量额外装入四倍空间分片。较大容量的维护影响在组织对照中另算。

每次选16个输出，八平面各配16差分ADC，共128个；一对差分支路同时进入一个转换器，不重复算两次标量转换。16个数字重构通道每轮两拍，输入逐bit、四输出组串行。128个pair写驱动含256支路控制，700 nA幅值在本地固定；一个外部更新域、128-bit编码接口。所有三情景的ADC、驱动和译码数量及额定能力固定。

令端点符号 $s=2b-1$，输入bit为$x$，64项差分和 $d=\sum xs$ 处于$-64\ldots64$，共129级；令共享输入popcount为$p=\sum x$，理想二进制部分和为$q=(d+p)/2$。实际名义服务先完成ADC量化再重构。采用固定的10-bit有符号码$c\in[-512,511]$，零电压码为0，LSB为$448/1024$ mV，向最近码量化、半码向正方向取整并在码端点裁剪。每单位$d$对应3.5 mV，因而连续仿射估计为$(c+8p)/16$。第一重构拍实施普通整数舍入
\[
\widehat q=\left\lfloor\frac{c+8p+8}{16}\right\rfloor.
\]
$c$符号扩展、7-bit的$p$左移三位、常数8并入既有偏移加法，右移四位得到整数部分和。名义0--64需7个无符号bit；运算保持有符号，以容纳任意ADC码与popcount组合给出的$-32\ldots64$，不隐藏裁剪。至少12-bit有符号临时值可容纳于既有22-bit重构数据路；16个通道各自已有八平面的并行前处理，第二拍保留位权累加，补码最高位取负权。不增加parity查表、分数累加器或额外求值轮次。此为明确的数字参考实现，未作门级时序验证。名义量化诊断同时保留整数舍入前的仿射残差与正满幅裁剪；22-bit输出容器、10-bit名义码宽和约8 ENOB分别记录，不将名义无噪声结果解释为模拟精度保证。

\section{原始证据与读写预算}

\begin{table}[htbp]\centering\small
\begin{tabularx}{\textwidth}{@{}>{\raggedright\arraybackslash}p{30mm}>{\raggedright\arraybackslash}X>{\raggedright\arraybackslash}X@{}}\toprule
来源定位 & 原值与条件 & 工程采用\\\midrule
GC-04 pp.6--7 Fig.10 & 64×64 pair、64个5-bit SAR、完整precharge/DTC/MAC/ADC周期180 ns。 & 180 ns仅作完整周期量级；本例二元输入与公共SAR前端独立列预算。\\
GC-04 pp.7--8 Figs.11--12 & 写BL约50 fF；5 ns电压粗写，随后电流精写；60 ns图为模拟。 & 保留两步写、动态cascode与负载，不能只取5 ns。\\
GC-04 pp.9--10 Fig.16 & 实测总65 ns传输函数达到目标，75 ns为较长测试点；均含5 ns粗写。 & 短/参考65 ns、长75 ns完整编程，不再额外叠粗写。\\
GC-04 p.10 Figs.17--18 & 端点±7，调RWL脉冲逐行ADC读回；400\,$\mu$s内99.7\%样本漂移小于原ADC 1 LSB。 & 固定400\,$\mu$s刷新；二元符号判决使用更大端点裕量，保留噪声/尾部资格条件。\\\bottomrule
\end{tabularx}
\caption{器件模式与原值保留。原ADC-LSB漂移统计不等于全部bit无误保证。}
\end{table}

读前端区分积分脉冲、共同阶段预算和控制窗口。二元MAC脉冲 $t_{\rm MAC}=1$ ns，单pair刷新符号读取脉冲 $t_{\rm RF}=64$ ns；它们均为本参考配置的工程选择。两种模式保留同一64-cell局部位线、700 nA端点、有效积分负载和ADC范围，但活动行数与RWL宽度不同。GC-04 p.10明确逐行保持读会调整RWL脉冲，支持操作模式的区别；原文没有报告本例1/64 ns或各阶段独立建立时间。

每支路有效积分负载200 fF。64行MAC的最大支路电荷和单pair刷新电荷均为 $64\times700\,\mathrm{nA}\times1\,\mathrm{ns}=44.8$ fC，对应最大224 mV；差分ADC覆盖$\pm224$ mV，128路双支路合计51.2 pF。相同满幅电荷仅核对量程，不证明两种模式的建立、漏电或噪声相同。名义10-bit转换、约8 ENOB和最终22-bit输出容器各有不同含义；容器不补回模拟有效分辨能力。

预充、局部选择、写隔离开关回接后的建立、隔离采样及恢复合并为共同预算 $t_c=36/86/136$ ns，保留所有必要阶段，不虚构各阶段的独立纳秒数。公共 $T_I$承担输入驱动，$T_A$承担转换，$T_D$承担已有数字控制和重构。主调度按当前模式的积分脉冲推进既有SAMP与SAR阶段：
\[
C_{A,M}=T_I+t_c+t_{\rm MAC}+T_A=49/112/197\ \mathrm{ns},\qquad
C_{A,F}=T_I+t_c+t_{\rm RF}+T_A=112/175/260\ \mathrm{ns}.
\]
GC-04 p.6--7的CCTRL、PRE、SAMP及可调RWL把预充、积分与已采样SAR分开；p.10逐行读取可调整RWL。二元输入每轮的所有活动RWL均在1 ns结束，不存在同轮未结束的多位输入脉冲，也没有必须等待另一个刷新操作64 ns脉冲的共享资源依赖。因此采用已有相位控制推进采样，无需新增ADC、驱动、样本保持阵列或shadow。1 ns后仍支付$t_c$内的隔离、建立和采样阶段；按模式推进不等于零后建立时间。

这两组前端时隙是相同资源和优化机会下的工程预算。模拟脉冲与相位边界采用已有可调延迟/控制，不强制每个模拟阶段为$T_D$整数拍；全部数字拍仍独立支付。原180 ns完整周期只锚定总量级，既不证明36/86/136 ns的分段建立，也不代表所算提速已被测量。将MAC采样边界保守地固定在64 ns的另一控制策略在后文单列。

写时隔离新增积分/SAR负载，将写BL含开关寄生保持在原约50 fF量级，才沿用65/75 ns完整建立锚点。原粗写约0.52 V、5 ns，在50 fF下每支路线性电容电流量级5.2\,$\mu$A；驱动、replica电流源和电压分配须支撑全部选中支路。128个pair最终目标电流合计89.6\,$\mu$A不等于粗写瞬态总供流。两相结构及新增隔离负载资格分别保留，不以裸电流目标代替完整驱动能力。

\section{完整求值、更新和无shadow刷新}

显式 $K=N=64$ 代入公共API，输入八片、一个64行组、四输出组：
\[
N_E=N_C=N_{\rm rec}=32,\quad N_{\rm dig}=64,\quad N_{\rm scalar}=4096,
\qquad C_S=32C_{A,M}+66T_D.
\]
典型 $C_S=32\times112+66\times5=3914$ ns。捕获与结果提交各一拍；阶段顺序执行，无跨向量流水或跨维护窗口模拟保持。

一笔resident完成16 Byte。128 pair需256支路选择bit，经128-bit口两数据拍，首拍随命令，准备/完成和额外数据拍共 $3T_D$；完整程序 $P=65/65/75$ ns含粗写和精写：
\[
C_R=3T_D+P=71/80/105\ \mathrm{ns},\qquad B_R=16\ \mathrm{Byte}.
\]
每pair两支路建立端点，电流反馈即编程机制，不另加未报告的program/verify或重试。256笔完整事务覆盖4096 Byte矩阵，原始完整装载 $T_{R,0}=256C_R$。

刷新逐逻辑行、逐16输出组扫描。128个差分ADC读取当前端点符号，128个本地sign decoder用一拍判决，暂存128bit符号，再编码成256支路控制bit并执行同样完整写。只有本组码和写控制保持，不保存全矩阵数字权重。每组
\[
C_F=C_{A,F}+T_D+C_R,\qquad N_F=64\times4=256,\qquad H=256C_F.
\]
一轮真实读取32768 pair、重写65536 cell；典型 $C_F=175+5+80=260$ ns、$H=66560$ ns。刷新不产生workload更新payload。原文64×65 ns=4.16\,$\mu$s只描述原宏逐行重写，没有覆盖本例无shadow读码流程，不能套用其1.1\%维护比例。

\section{周期调度、三情景与两表接口}

每400\,$\mu$s按固定起点、固定顺序刷新；同一cell的写回位置每周期相同。刷新开始前 $G$ 禁止发起新不可抢占组，其中
\[
G=\max(C_{A,M}+2T_D,C_R)=71/122/217\ \mathrm{ns}.
\]
计算在完成一个求值/转换/重构组后暂停，数字累加器和输入寄存器继续保持；刷新符号暂存128 bit及写控制256 bit与它们独立，ADC码在重构完成后可被刷新复用。写在完整16 Byte组边界暂停。单次guard覆盖可能空隙，无跨窗口模拟样本，也不需要额外全矩阵shadow。冷装载后建立刷新相位；正常情景每周期至少303.783\,$\mu$s有用窗口，远大于217 ns最大组，足以安排这些分批服务。

当两路分别使用同一共享资源时，固定预留给出
\[
\alpha=1-(H+G)/400000,\qquad
\rho=64\alpha/C_S,\quad\tau=16\alpha/C_R.
\]
$\alpha\le0$时调度不可行，有效时间和能力为空。正可用率仍依赖上述期限、批边界和寄存器生命期；服务间隔是稳态平均，单请求latency取决于到达相位。外部新写不自动抵扣其他cell维护；可利用完整新写抵扣的另一调度需重算，不混入固定预留能力。

% Generated by scripts/check_gain_cell_edram.py --emit.
\begin{table}[htbp]\centering\small
\begin{tabular}{@{}lrrrrrr@{}}\toprule
情景 & $C_A$ (ns) & $P$ (ns) & $C_S$ (ns) & $C_R$ (ns) & $T_{\rm ref}$ ($\mu$s) & 可用率\\\midrule
短 & 112 & 65 & 3716 & 71 & 47.360 & 88.131\%\\
参考 & 175 & 65 & 5930 & 80 & 66.560 & 83.314\%\\
长 & 260 & 75 & 8980 & 105 & 96.000 & 75.930\%\\
\bottomrule\end{tabular}
\caption{无维护时的占用 $C_S,C_R$ 与每 0.4 ms 全矩阵刷新时间；可用率再扣批边界余量 116/185/280 ns。}\end{table}
\begin{table}[htbp]\centering\small
\begin{tabular}{@{}lrrrrr@{}}\toprule
情景 & $\Delta_S$ ($\mu$s) & $\Delta_R$ (ns) & $\rho$ (MB/s) & $\tau$ (MB/s) & $\RI^*$\\\midrule
短 & 4.216 & 80.56 & 15.18 & 198.6 & 0.07643\\
参考 & 7.118 & 96.02 & 8.992 & 166.6 & 0.05396\\
长 & 11.827 & 138.29 & 5.411 & 115.7 & 0.04677\\
\bottomrule\end{tabular}
\caption{包含固定刷新预留后的长期平均服务；$B_S=64$ Byte、$B_R=16$ Byte，MB=$10^6$ Byte。三点均为同一64×64原生组织的可持续情景；容量压力另列。}\label{09_gain_cell_edram:tab:results}\end{table}


典型 $\alpha=(400000-66560-122)/400000=0.833295$；原始能力为16.3516/200 MB/s，维护后为13.6257/166.659 MB/s，$\RI^*\approx0.08176$。原始装载20480 ns，按同一周期预留的周期长期完整装载平均服务占用为 $T_R=20480/\alpha\approx24577.1$ ns，不是单次冷装载latency；相容的完整向量平均间隔为 $3914/\alpha\approx4697.02$ ns。因此
\[
U^*=T_R/\Delta_S=20480/3914=64\RI^*\approx5.23250.
\]
每16 KiB展示成本为此4 KiB原生装载平均的四倍，不能称一次真实16 KiB请求延迟。

典型读前端3584 ns占原始计算约92\%，写编程65 ns占80 ns事务约81\%；全矩阵读码重写的维护损失约17\%。原生64×64的短/参考/长可用率约88\%/83\%/76\%，均有持续服务余量；维护临界点另列。

\section{保守控制对照及共享参数依赖}

固定控制策略让两种读取均到64 ns预留窗口结束后才推进采样。MAC脉冲仍只积分1 ns，余下63 ns要求模拟状态保持，不能解释为器件必需积分或建立。该保守对照维持同一$t_c$、偏置、负载、量程、写入和保持期；求值前端变为112/175/260 ns，刷新前端及全刷新占用不变。短/参考/长的原始$C_S$为3716/5930/8980 ns，guard为116/185/280 ns，长期$\rho$为15.1786/8.99170/5.41149 MB/s，$\tau$为198.605/166.628/115.703 MB/s。

典型主调度相对固定窗口的输入能力比为1.51536，来自32个计算组中未用的控制预留；更新能力的小幅差别仅来自guard。短主调度的guard由完整写事务决定：$G=\max(49+4,71)=71$ ns，故不能把每种场景或每个前端都减63 ns。两种调度均保留逐pair读回、符号译码、完整两步写及无shadow维护。

共同$T_I,T_A,T_D,t_c$绑定实际消费者；改变同一参数会重算计算和刷新，继而影响$H,G,\alpha$和两路有效能力。主模式的MAC脉冲改变正常计算与guard，刷新脉冲改变全刷新；完整写程序同时用于resident更新与刷新写回。固定窗口内改变积分脉冲不改变预留占用，但不允许脉冲超过控制槽。控制对照不代替共同前端参数的不确定性传播。

检查器以32次计算、256次刷新、每次写3拍和符号译码1拍独立计数，核对九项参数扰动；另核对宽32输出、较大容量、短情景写主导guard、临界和不可行情景以及完整装载$U^*$。共同无噪声诊断采用声明的理想积分、固定量化器及有限整数重构，并保存舍入前损失。这些检查验证模型语义、参数传播与算术，不代替电路时序、噪声或重复刷新可靠性验证。

\section{容量及资源对照与适用范围}

保持局部64行负载、128 ADC和128 pair驱动，扩成K128/N128需32个tile、四倍容量；对应1024刷新组。典型原始计算15626 ns、刷新266240 ns、guard122 ns、可用率33.4095\%，维护后约2.737/66.819 MB/s。该容量对照不是主三情景，增加容量而不增加共享前端会显著增加每周期维护负担。

大容量长条件的刷新384000 ns，加guard217 ns，可用率仅3.94575\%，单列临界压力。令刷新专用前端$t_f=t_c+64$ ns，同一$t_c$也作用正常MAC；其他长参数固定时须满足 $1025t_f+179217<400000$ ns，即$t_f<215.39805$ ns。$t_f=216$ ns时刷新占400384 ns、guard233 ns，总占400617 ns，平均能力和RI*为空。没有延长400\,$\mu$s保持期或裁剪负可用率。

原生64×64的32输出资源对照将差分ADC、pair驱动和符号decoder均翻倍至256路，数字重构保持16通道；原始求值2122 ns，32 Byte写90 ns，刷新34560 ns，guard132 ns。所有资源增加显式列账：
\begin{table}[htbp]\centering\small
\begin{tabular}{@{}lrrrrrr@{}}\toprule
配置 & ADC/写对驱动 & $N_E$ & 刷新 ($\mu$s) & $\rho$ (MB/s) & $\tau$ (MB/s) & $\RI^*$\\\midrule
16 输出 & 128/128 & 32 & 66.56 & 13.63 & 166.7 & 0.08176\\
32 输出 & 256/256 & 16 & 34.56 & 27.54 & 324.7 & 0.08483\\
\bottomrule\end{tabular}
\caption{仅参考时隙的资源对照；32 输出配置每事务完成 32 Byte，数据接口仍为 128 bit、数字通道仍为 16。}\label{09_gain_cell_edram:tab:wide}\end{table}


\medskip
\noindent\textbf{核心来源。} J. Song et al., JSSC 2024，\href{https://doi.org/10.1109/JSSC.2023.3339887}{GC-04} pp.6--10，Figs.8、10--12、16--18。完整来源、实现条件及复算账本见本例输入JSON与证据笔记。
